\scanpage{41}
\mspage{8}
J'avais craint un moment, par cet exemple,
que le type total d'homotopie \inserted{de $X$} ne pourrait
pas (avec les hypothèses \inserted{$K(\pi,1)$} faites) se récupérer
à partir de ceux de $X^*$, $Y$, $\mathcal V^*_{Y,X}$ et du
diagramme (7), i.e du diagramme (8)
ou (10). Mais je me suis convaincu ensuite
que cette crainte n'était pas fondée, $\unsure{\overline m|=}$ si
le topos $\unsure{\mathcal F|\mathfrak F}$ n'est pas suffisant, lui, pour
exprimer le type d'homotopie \unsure{correct|courant}.
Essayons en effet, dans le cas qui nous
occupe --- le cas de la division cellulaire
la plus simple de $S^{\unsure{2|7}}$, (qui est en $\unsure{=|\overline m}$ temps
\deleted{élémentaire}
le cas le plus \deleted{simple} élémentaire d'un
diviseur sur une courbe alg complexe
compacte !) --- de récupérer le $H^2(X,A)$
($A$ un groupe de coeff comm, disons),
à partir du diagramme (23). J'interprète
$H^2(X,A)$ comme classifiant les \uline{gerbes} sur
$X$ liées par $A$ --- et je note que pour des
\deleted{telles} gerbes \inserted{sur des topos}, on a une théorie de ``recolle-
% Pale bottom-left foliation, another hand; distinct from the author's 8.
\margin{40}
